\subsection{বায়োমেট্রিক্স}

\begin{learningbox}
এই অংশ শেষে শিক্ষার্থী পারবে:
\begin{itemize}
\item বায়োমেট্রিক্সের ধারণা ব্যাখ্যা করতে।
\item fingerprint, face, iris, retina ও voice recognition-এর ব্যবহার লিখতে।
\item authentication ও identification-এর পার্থক্য বুঝতে।
\item বায়োমেট্রিক্সের সুবিধা, সীমাবদ্ধতা ও privacy risk বিশ্লেষণ করতে।
\end{itemize}
\end{learningbox}

\subsubsection{বায়োমেট্রিক্সের ধারণা}

মানুষের কিছু শারীরিক ও আচরণগত বৈশিষ্ট্য তুলনামূলকভাবে অনন্য। যেমন আঙুলের ছাপ, মুখমণ্ডল, চোখের iris, কণ্ঠস্বর বা হাতের লেখা। এসব বৈশিষ্ট্য digital পদ্ধতিতে মেপে পরিচয় যাচাই বা শনাক্ত করার প্রযুক্তিকে বায়োমেট্রিক্স বলা হয়।

\begin{definitionbox}{সংজ্ঞা}
বায়োমেট্রিক্স হলো মানুষের অনন্য শারীরিক বা আচরণগত বৈশিষ্ট্য digitalভাবে সংগ্রহ, বিশ্লেষণ ও তুলনা করে পরিচয় যাচাই বা শনাক্ত করার প্রযুক্তি।
\end{definitionbox}

NIST-এর biometric guidance অনুযায়ী biometric system সাধারণত capture, feature extraction, template storage এবং matching ধাপে কাজ করে। অর্থাৎ সরাসরি “আঙুলের ছবি” নয়; অনেক ক্ষেত্রে বৈশিষ্ট্য থেকে template তৈরি করে তুলনা করা হয়।

\subsubsection{ফিঙ্গারপ্রিন্ট শনাক্তকরণ}

Fingerprint recognition আঙুলের ridge, valley, minutiae point ইত্যাদি বিশ্লেষণ করে। Mobile unlock, attendance system, passport/ID verification এবং criminal investigation-এ fingerprint ব্যবহৃত হয়।

\subsubsection{মুখমণ্ডল শনাক্তকরণ}

Face recognition camera image থেকে মুখের feature বিশ্লেষণ করে template তৈরি করে। Smartphone unlock, airport security, surveillance এবং photo organization-এ ব্যবহৃত হয়। আলো, angle, mask বা image quality performance কমাতে পারে।

\subsubsection{আইরিস ও রেটিনা স্ক্যান}

Iris হলো চোখের রঙিন অংশ; এর pattern খুব অনন্য। Retina scan চোখের পেছনের blood vessel pattern বিশ্লেষণ করে। উচ্চ নিরাপত্তা দরকার এমন ক্ষেত্রে iris/retina-based system ব্যবহার করা যায়।

\subsubsection{ভয়েস রিকগনিশন}

Voice recognition কণ্ঠস্বরের pitch, tone, rhythm ও speech pattern বিশ্লেষণ করে। Call center verification, voice assistant এবং accessibility system-এ ব্যবহার হয়। তবে noise, অসুস্থতা বা mimicry problem হতে পারে।

\subsubsection{বায়োমেট্রিক্সের ব্যবহারক্ষেত্র}

\begin{keypointbox}{ব্যবহারক্ষেত্র}
\begin{itemize}
\item জাতীয় পরিচয়পত্র, পাসপোর্ট ও immigration।
\item অফিস attendance ও access control।
\item smartphone unlock ও mobile payment।
\item ব্যাংকিং ও financial verification।
\item আইনশৃঙ্খলা ও forensic investigation।
\item পরীক্ষাকেন্দ্র বা secure lab-এ identity verification।
\end{itemize}
\end{keypointbox}

\begin{center}
\begin{tabular}{|p{0.30\textwidth}|p{0.56\textwidth}|}
\hline
\textbf{ধরন} & \textbf{উদাহরণ} \\
\hline
শারীরিক বৈশিষ্ট্য & fingerprint, face, iris, retina, palm vein \\
\hline
আচরণগত বৈশিষ্ট্য & voice, signature, typing rhythm, gait \\
\hline
\end{tabular}
\end{center}

\begin{keypointbox}{সুবিধা ও সীমাবদ্ধতা}
\begin{itemize}
\item সুবিধা: password মনে রাখতে হয় না, পরিচয় যাচাই দ্রুত, নকল করা কঠিন।
\item সীমাবদ্ধতা: sensor error, false match, injury বা lighting problem, privacy risk।
\item নিরাপত্তা: biometric data leak হলে password-এর মতো সহজে বদলানো যায় না।
\end{itemize}
\end{keypointbox}

\begin{boardbox}{বোর্ড ফোকাস}
বায়োমেট্রিক্সের উত্তরে “মানুষের অনন্য বৈশিষ্ট্য” শব্দটি অবশ্যই রাখবে। fingerprint, face, iris, voice—কমপক্ষে চারটি উদাহরণ মনে রাখো।
\end{boardbox}

\begin{practicebox}
\textbf{MCQ}
\begin{enumerate}
\item আঙুলের ছাপ দিয়ে পরিচয় যাচাই কোন প্রযুক্তির উদাহরণ?\\
ক. Biometrics \quad খ. Blogging \quad গ. Formatting \quad ঘ. Spreadsheet
\item চোখের রঙিন অংশের pattern শনাক্তকরণ কোনটি?\\
ক. Iris recognition \quad খ. OCR \quad গ. GPS \quad ঘ. Router
\end{enumerate}

\textbf{সংক্ষিপ্ত প্রশ্ন}
\begin{enumerate}
\item বায়োমেট্রিক্স কী?
\item বায়োমেট্রিক্সের দুটি privacy risk লেখ।
\end{enumerate}
\end{practicebox}

\begin{sourcebox}
এই অংশে local ICT PDF resources, NIST biometric concepts, ISO-style biometric vocabulary overview এবং HSC ICT resources মিলিয়ে লেখা হয়েছে।
\end{sourcebox}
